%% Question bank chapter file - generated by bank_tools.py build (do not edit while the bank is still being built:
%% it is regenerated from the per-test data in Question Bank\data). Every question is one bankq block:
%%   \begin{bankq}[id=..., source={...}, paper=.., qno=.., type=..., subtopic={...}, answer={...}, answer_src=key|solved,
%%                 working={...}, difficulty=1-10, diff_parts={K S I C F}, diff_note={...},
%%                 group=<passage id or empty>, also_in={duplicate copies in other tests}]
%% difficulty = min(10, K+S+I+C+F-1): K concepts 1-3, S steps 1-4, I insight 0-2, C computation 0-2, F format 0-1
%%   ...question LaTeX, options, figures (\bfig / \bfigi -> img/)...
%%   \end{bankq}
%% Passages shared by several questions are bankpassage blocks; their questions carry group=<passage id>.
\documentclass[10pt]{article}
\usepackage[a4paper,margin=1.8cm]{geometry}
\usepackage[T1]{fontenc}
\usepackage{amsmath,amssymb,mathtools}
\usepackage{graphicx}
\usepackage[export]{adjustbox}
\usepackage{xcolor,needspace,keyval}
\usepackage{enumitem}
\usepackage{booktabs,array,tabularx}
\usepackage[hidelinks]{hyperref}
\setlength{\parindent}{0pt}
\setlength{\parskip}{3pt}
\makeatletter
\define@key{bq}{id}{\def\bq@id{#1}}\define@key{bq}{source}{\def\bq@src{#1}}\define@key{bq}{paper}{\def\bq@paper{#1}}
\define@key{bq}{qno}{\def\bq@qno{#1}}\define@key{bq}{type}{\def\bq@type{#1}}\define@key{bq}{subtopic}{\def\bq@sub{#1}}
\define@key{bq}{answer}{\def\bq@ans{#1}}\define@key{bq}{answer_src}{\def\bq@asrc{#1}}\define@key{bq}{working}{\def\bq@work{#1}}
\define@key{bq}{group}{\def\bq@group{#1}}\define@key{bq}{also_in}{\def\bq@also{#1}}
\define@key{bq}{difficulty}{\def\bq@diff{#1}}\define@key{bq}{diff_parts}{\def\bq@dparts{#1}}\define@key{bq}{diff_note}{\def\bq@dnote{#1}}
\newenvironment{bankq}[1][]{\def\bq@work{}\def\bq@also{}\def\bq@group{}\def\bq@diff{?}\def\bq@dparts{}\def\bq@dnote{}\setkeys{bq}{#1}%
  \par\needspace{4\baselineskip}\medskip\noindent{\small\textbf{\bq@id}\hfill\color{gray}[\bq@src, Paper \bq@paper, Q\bq@qno]
  \bq@type\ $\cdot$ \bq@sub\ $\cdot$ \textbf{D\bq@diff}}\par\nopagebreak\relax}
  {\par{\small\color{teal}Answer: \textbf{\bq@ans}\ifx\bq@work\@empty\else\ (solved: \bq@work)\fi
  \ifx\bq@also\@empty\else\par\color{gray}Also in: \bq@also\fi}\par}
\define@key{bp}{id}{\def\bp@id{#1}}\define@key{bp}{source}{\def\bp@src{#1}}\define@key{bp}{questions}{\def\bp@qs{#1}}
\newenvironment{bankpassage}[1][]{\setkeys{bp}{#1}\par\needspace{4\baselineskip}\bigskip\noindent
  {\small\textbf{Passage \bp@id}\hfill\color{gray}[\bp@src; questions \bp@qs]}\par\nopagebreak\relax\itshape}{\par}
\makeatother
\newcommand{\bfig}[2]{\par\smallskip{\centering\includegraphics[width=#1pt,max width=0.62\linewidth,max height=0.32\textheight]{img/#2}\par}\smallskip}
\newcommand{\bfigi}[2]{\includegraphics[width=#1pt,max width=\linewidth,max height=0.32\textheight,valign=c]{img/#2}}
\newcommand{\bmissing}[1]{\textit{[figure missing in the source conversion: \detokenize{#1}]}}
\newenvironment{qoptions}{\begin{itemize}[leftmargin=1.8em,itemsep=1pt,topsep=2pt,label={}]}{\end{itemize}}
\newcommand{\opt}[2]{\item[(#1)] #2}
\providecommand{\fwstart}{}\providecommand{\fwend}{}\providecommand{\qfullsplit}{}
\newcolumntype{L}[1]{>{\raggedright\arraybackslash}p{#1}}
\newcolumntype{C}[1]{>{\centering\arraybackslash}p{#1}}
\newcolumntype{R}[1]{>{\raggedleft\arraybackslash}p{#1}}
\providecommand{\sectionhdr}[1]{}\providecommand{\subsectionhdr}[1]{}


\begin{document}
{\Large\bfseries Gravitation --- Question Bank}\par\smallskip
{\small Physics chapter P06. Questions from all converted test papers, grouped by source series. Answers: printed key where it matches the paper, otherwise solved (answer\_src=solved, with working). Exact repeats across papers appear once, with also\_in listing the other copies.}\par\medskip\hrule

\section*{Score-I papers}
\begin{bankq}[id=SI12-P1-PHY-S3-Q06, source={Score-I - Test 12 (09 Mar 2022)}, paper=1, qno=6, type=numerical, subtopic={Field and potential}, answer={4}, answer_src=key, difficulty=2, diff_parts={K1 S1 I0 C0 F1}, diff_note={concepts: variation of g with height | steps: g' = g(R/(R+h))\^{}2 = 0.64 g gives h = R/4}]
The height at which the acceleration due to gravity decreases by 36\% of its value on the surface of the earth is $R/x$, then the value of $x$ is: (Assume radius of the earth is $R$)
\end{bankq}

\section*{Score Advanced papers}
\begin{bankq}[id=SA01-P1-PHY-S1-Q02, source={Score Advanced - Test 01 (22 Sep 2020)}, paper=1, qno=2, type=single, subtopic={Field and potential}, answer={C}, answer_src=key, difficulty=6, diff_parts={K2 S2 I2 C1 F0}, diff_note={concepts: two bodies under mutual attraction, field and potential | steps: momentum conservation; energy conservation; speeds and time}]
Two small bodies of masses 1 kg and 4 kg are at rest on a smooth horizontal surface. When they are at a separation of 10 m, they are released to move under mutual attractive interaction. Then, which of the following statements is \textbf{CORRECT}?\par
\bfig{89.9}{SA01_p1_phy_s1_q02_1.png}
\begin{qoptions}
\opt{A}{They will never meet}
\opt{B}{They will meet at a distance of 6 m from A.}
\opt{C}{To an observer at the centre of mass, the particles have equal and opposite momenta.}
\opt{D}{They will meet at a distance of 8 m from B.}
\end{qoptions}
\end{bankq}
\begin{bankq}[id=SA01-P1-PHY-S1-Q06, source={Score Advanced - Test 01 (22 Sep 2020)}, paper=1, qno=6, type=single, subtopic={Kepler's laws}, answer={C}, answer_src=key, difficulty=6, diff_parts={K2 S2 I2 C1 F0}, diff_note={concepts: three planets in a line, kepler's laws | steps: periods from Kepler law; alignment condition; time}]
There are three planets in circular orbits around a fixed star at distances a, 4a and 9a, respectively in same sense (say anticlockise). At time t = 0, the star and the three planets are in a straight line. The period of revolution of the closest planet is T. What minimum time after t = 0 will they again be in the same straight line?
\begin{qoptions}
\opt{A}{8 T}
\opt{B}{27 T}
\opt{C}{108 T}
\opt{D}{512 T}
\end{qoptions}
\end{bankq}
\begin{bankq}[id=SA01-P2-PHY-S2-Q09, source={Score Advanced - Test 01 (22 Sep 2020)}, paper=2, qno=9, type=single, subtopic={Satellites}, answer={A}, answer_src=key, difficulty=7, diff_parts={K2 S2 I2 C2 F0}, diff_note={concepts: particle given a tangential velocity at distance r, satellites | steps: energy; angular momentum; orbit type; match}]
A particle is taken to a distance r ( $>$ R) from centre of earth. R is the radius of earth. It is given velocity v which is perpendicular to $\vec r$ (measured from centre of earth). Match the following assuming particle do not collide with earth. Mass of particle (m) is negligible in comparision to mass of earth (M) and assume there is no air resistance :\par
\noindent{\footnotesize\setlength{\tabcolsep}{2pt}\begin{tabular}{|L{.8cm}|L{3.0cm}|L{.8cm}|L{3.2cm}|}\hline
\multicolumn{2}{|c|}{List--I} & \multicolumn{2}{c|}{List--II}\\\hline
(P) & $v=\sqrt{\dfrac{GM}r}$ & (1) & Distance of point mass from centre of earth first increases then decreases\\\hline
(Q) & $v=\sqrt{\dfrac{2GM}r}$ & (2) & Distance of point mass from centre of earth first decreases then increases\\\hline
(R) & $\sqrt{\dfrac{2GM}r}>v>\sqrt{\dfrac{GM}r}$ & (3) & Distance of point mass from centre of earth monotonically increases\\\hline
(S) & $\sqrt{\dfrac{GM}r}>v>\sqrt{\dfrac{GM}r\left(\dfrac{2R}{r+R}\right)}$ & (4) & Distance of point mass from centre of earth remain constant\\\hline
\end{tabular}}\par
\begin{qoptions}
\opt{A}{P $\to$ 4 ; Q $\to$ 3 ; R $\to$ 1 ; S $\to$ 2}
\opt{B}{P $\to$ 3 ; Q $\to$ 4 ; R $\to$ 2 ; S $\to$ 1}
\opt{C}{P $\to$ 1 ; Q $\to$ 2 ; R $\to$ 3 ; S $\to$ 4}
\opt{D}{P $\to$ 4 ; Q $\to$ 3 ; R $\to$ 2 ; S $\to$ 1}
\end{qoptions}
\end{bankq}
\begin{bankq}[id=SA01-P2-PHY-S2-Q12, source={Score Advanced - Test 01 (22 Sep 2020)}, paper=2, qno=12, type=single, subtopic={Field and potential}, answer={D}, answer_src=key, difficulty=7, diff_parts={K2 S2 I2 C2 F0}, diff_note={concepts: particle falling on a box, field and potential | steps: relative acceleration; fall time; motion of the box}]
Figure shows a box B of mass m held fixed at large height H (which is very large in comparison to dimension of the box, but very small in comparison to radius of earth). A particle P whose mass is equal to that of box is suspended at O (which is centre of mass of the box) with a massless string. Particle can swing in vertical plane without colliding box. Particle is pulled aside to horizontal position and released. The box is held fixed till the particle reaches lowest position and then box is released at t = 0. t is a general time before box collides ground.\par
\bfig{75.3}{SA01_p2_phy_s2_q12_1.png}
\noindent{\footnotesize\setlength{\tabcolsep}{2pt}\begin{tabular}{|L{.8cm}|L{4.8cm}|L{.8cm}|L{1.4cm}|}\hline
\multicolumn{2}{|c|}{List-I} & \multicolumn{2}{c|}{List-II}\\\hline
(P) & Tension in the thread at the t = t in the units of mg will be & (1) & 1\\\hline
(Q) & Speed of box at time t = $\pi\sqrt{\dfrac\ell{2g}}$ in the units of $\sqrt{\dfrac{g\ell}8}\sqrt{\pi^2+4}$ will be & (2) & 2\\\hline
(R) & Acceleration of the box at t = $\pi\sqrt{\dfrac\ell{2g}}$ in the units of g will be & (3) & 0\\\hline
(S) & Initial acceleration (t = 0$^+$) of particle in the units of g will be & (4) & none\\\hline
\end{tabular}}\par
\begin{qoptions}
\opt{A}{P $\to$ 3 ; Q $\to$ 2 ; R $\to$ 1 ; S $\to$ 4}
\opt{B}{P $\to$ 1 ; Q $\to$ 1 ; R $\to$ 3 ; S $\to$ 4}
\opt{C}{P $\to$ 4 ; Q $\to$ 1 ; R $\to$ 3 ; S $\to$ 2}
\opt{D}{P $\to$ 1 ; Q $\to$ 2 ; R $\to$ 3 ; S $\to$ 3}
\end{qoptions}
\end{bankq}
\begin{bankq}[id=SA04-P1-PHY-S1-Q03, source={Score Advanced - Test 04 (04 Jul 2021)}, paper=1, qno=3, type=multi, subtopic={Field and potential}, answer={A,C}, answer_src=solved, working={The cylinder field exceeds the sphere field for d much greater than R, so the field and the orbital speed both increase}, difficulty=7, diff_parts={K2 S2 I2 C1 F1}, diff_note={concepts: gravitational field of a long cylinder versus a sphere, field and potential | steps: field of an infinite cylinder 2G(lambda)/d; field of a sphere GM/d\^{}2; ratio (3/2)(d/R) is large; orbital speed v\^{}2=g d}]
Assume that the earth changes its shape and turns into an infinite cylinder whose radius and distance of moon from the axis of cylindrical earth remains unchanged and also moon is remain spherical. Assume density of earth remains same and distance of moon from earth is much greater than radius of earth. Which of the following statement(s) is/are correct :
\begin{qoptions}
\opt{A}{Gravitational field at the moon due to earth increases.}
\opt{B}{Gravitational field at the moon due to earth decreases.}
\opt{C}{Speed of the moon in its orbit around the cylindrical earth increases.}
\opt{D}{Speed of the moon in its orbit around the cylindrical earth decreases.}
\end{qoptions}
\end{bankq}
\begin{bankq}[id=SA04-P1-PHY-S1-Q04, source={Score Advanced - Test 04 (04 Jul 2021)}, paper=1, qno=4, type=multi, subtopic={Satellites}, answer={B,C}, answer_src=solved, working={Perihelion: a(1-e)=3R/2 gives a=3R, E=-GMm/6R. Aphelion: a(1+e)=3R/2 gives a=R, E=-GMm/2R. Changes GMm/6R and GMm/2R}, difficulty=7, diff_parts={K2 S2 I2 C1 F1}, diff_note={concepts: energy change between a parabolic path and ellipses, satellites | steps: parabola has E=0 with closest distance 3R/2; perihelion case a=3R; aphelion case a=R; E=-GMm/2a}]
A particle of mass m is thrown in a parabolic path towards earth having mass M and radius R. It comes to a greatest approach distance R/2 from surface of earth. Magnitude of total energy change required to make the particle moving in an elliptical path having eccentricity 1/2, with the greatest approach point either aphelion or perihelion are :
\begin{qoptions}
\opt{A}{$\dfrac{1}{3}\dfrac{\text{GMm}}{R}$}
\opt{B}{$\dfrac{1}{6}\dfrac{\text{GMm}}{R}$}
\opt{C}{$\dfrac{1}{2}\dfrac{\text{GMm}}{R}$}
\opt{D}{$\dfrac{\text{GMm}}{R}$}
\end{qoptions}
\end{bankq}
\begin{bankq}[id=SA04-P1-PHY-S2-Q03, source={Score Advanced - Test 04 (04 Jul 2021)}, paper=1, qno=3, type=multi, subtopic={Satellites}, answer={B,C}, answer_src=solved, working={T=84 min, 17.2 revolutions per day, dr per revolution=(2/3)R(dT/T)=53 m=160/3 m. Angular momentum sqrt(GMr) decreases and the energy -GMm/2r decreases}, difficulty=7, diff_parts={K2 S2 I2 C1 F1}, diff_note={concepts: decay of a near-earth orbit, satellites | steps: T proportional to r\^{}1.5; fall per revolution; angular momentum and energy}]
A near earth orbit satellite of earth experiences very small air resistance. It's period of revolution decreases by 1.08 sec/day.
\begin{qoptions}
\opt{A}{It's radius decreases by $\dfrac{160}{3}$ m in 1 day.}
\opt{B}{It's radius decreases by $\dfrac{160}{3}$ m in 1 revolution.}
\opt{C}{It's angular momentum decreases}
\opt{D}{It's mechanical energy increases.}
\end{qoptions}
\end{bankq}
\begin{bankq}[id=SA04-P1-PHY-S2-Q04, source={Score Advanced - Test 04 (04 Jul 2021)}, paper=1, qno=4, type=multi, subtopic={Satellites}, answer={A,B,C,D}, answer_src=solved, working={VA/VB=2, EA/EB=(3/r)/(1/4r)=12, uA/uB=2, TA/TB=1/8}, difficulty=6, diff_parts={K2 S2 I1 C1 F1}, diff_note={concepts: ratios for two circular orbits, satellites | steps: V=sqrt(GM/r); E=-GMm/2r; escape speed sqrt(2GM/r); T proportional to r\^{}1.5}]
Satellite A is in a circular orbit of radius r and satellite B in another circular orbit of radius 4 r, both revolving around the earth. The masses of the satellites A and B are in the ratio 3 : 1. If the symbols V, E, u and T represent the speed, total energy, escape velocity and period of revolution of satellite, then
\begin{qoptions}
\opt{A}{V$_\text{A}$/V$_\text{B}$ = 2}
\opt{B}{E$_\text{A}$/E$_\text{B}$ = 12}
\opt{C}{u$_\text{A}$/u$_\text{B}$ = 2}
\opt{D}{T$_\text{A}$/T$_\text{B}$ = $\dfrac{1}{8}$}
\end{qoptions}
\end{bankq}
\begin{bankq}[id=SA04-P2-PHY-S1-Q15, source={Score Advanced - Test 04 (04 Jul 2021)}, paper=2, qno=15, type=single, subtopic={Satellites}, answer={D}, answer_src=solved, working={Same density as the earth gives the same period as an earth-surface orbit, 84 min}, difficulty=4, diff_parts={K2 S1 I1 C1 F0}, diff_note={concepts: orbital period of a pebble about an asteroid, satellites | steps: T=sqrt(3 pi/(G rho)) depends only on density}]
A spherical asteroid having the same density as that of earth is floating in free space. A small pebble is revolving around the asteroid under the influence of gravity near the surface of the asteroid. What is the approximate time period of the pebble?
\begin{qoptions}
\opt{A}{24 h}
\opt{B}{365 days}
\opt{C}{10 min}
\opt{D}{1 hr 24 min}
\end{qoptions}
\end{bankq}
\begin{bankq}[id=SA04-P2-PHY-S1-Q19, source={Score Advanced - Test 04 (04 Jul 2021)}, paper=2, qno=19, type=single, subtopic={Field and potential}, answer={D}, answer_src=solved, working={V0 R=v a and V0\^{}2/2 = v\^{}2/2 - GM/a give a=R(1+2GM/(aV0\^{}2))\^{}(-1/2)}, difficulty=6, diff_parts={K2 S2 I2 C1 F0}, diff_note={concepts: distance of closest approach to a black hole, field and potential | steps: angular momentum conservation; energy conservation; eliminate v}]
A mass m, travelling at speed V$_0$ in a straight line from far away is deflected when it passes near a black hole of mass M which is at a perpendicular distance R from the original line of flight. a, the distance of closest approach between the mass and the black hole is given by the relation :\par
\bfig{153.8}{SA04_p2_phy_s1_q19_1.png}
\begin{qoptions}
\opt{A}{$a=R\left(1+\dfrac{2GM}{aV_0^2}\right)^{1/2}$}
\opt{B}{$a=R\left(1+\dfrac{aV_0^2}{2GM}\right)^{1/2}$}
\opt{C}{$a=R\left(1+\dfrac{GM}{2aV_0^2}\right)^{-1/2}$}
\opt{D}{$a=R\left(1+\dfrac{2GM}{aV_0^2}\right)^{-1/2}$}
\end{qoptions}
\end{bankq}
\begin{bankq}[id=SA05-P1-PHY-S2-Q07, source={Score Advanced - Test 05 (15 Jul 2021)}, paper=1, qno=7, type=numerical, subtopic={Field and potential}, answer={2}, answer_src=key, difficulty=6, diff_parts={K2 S1 I2 C1 F1}, diff_note={concepts: speed through a tunnel through the earth, field and potential | steps: potential at the centre -3GM/2R; energy conservation}]
A tunnel is dug inside the earth across one of its diameters. A particle is projected inside the tunnel with velocity $\sqrt{\dfrac{\text{GM}}{\text{R}}}$ from one of its ends. Then maximum velocity attained by the particle in the subsequent motion is $\sqrt{\dfrac{\text{nGM}}{\text{R}}}$, then value of n. (assuming tunnel to be frictionless)
\end{bankq}
\begin{bankq}[id=SA07-P2-PHY-S3-Q06, source={Score Advanced - Test 07 (16 Apr 2023)}, paper=2, qno=6, type=numerical, subtopic={Escape velocity}, answer={3}, answer_src=key, difficulty=8, diff_parts={K3 S1 I2 C2 F1}, diff_note={concepts: speed to reach the second planet, escape velocity, point of zero net field | steps: point of zero net field; energy to cross it}]
Two planets of mass M \& 4M are shown in the figure. Radius of the planets R and 2R respectively and distance between centres of the planets in 6R. A small particle is thrown from planet of mass M along the line joining centres of both particles. Minimum speed of the particle so that it can reach other planet is given by $\sqrt{\dfrac{\text{x}\cdot\text{GM}}{5\text{R}}}$. Find x.\par
\bfig{176.5}{SA07_p2_phy_s3_q06_1.png}
\end{bankq}

\section*{Enthusiast Pre-Score papers}
\begin{bankq}[id=EPS02-P1-PHY-S2-Q08, source={Enthusiast Pre-Score - Test 02 (11 Nov 2020)}, paper=1, qno=8, type=multi, subtopic={Field and potential}, answer={A,C}, answer_src=key, difficulty=9, diff_parts={K2 S3 I2 C2 F1}, diff_note={concepts: particle sliding in a chord tunnel of the earth, field and potential | steps: force along the chord; SHM; amplitude from speed; time; maximum speed}]
A tunnel is dug across the earth along a chord AB as shown. A particle is projected inside the tunnel from end `A' with an initial velocity $v_0=\dfrac{\sqrt{3g_0R}}{2}$, there is no friction between the tunnel and the particle. Then choose the correct option(s). (where g$_0$ = acceleration due to gravity on the surface of Earth, R = radius of earth).\par
\bfig{138.7}{EPS02_p1_phy_s2_q08_1.png}
\begin{qoptions}
\opt{A}{The time taken by the particle to move from end `A' to other end `B' of the tunnel is $\dfrac{\pi}{3}\sqrt{\dfrac{R}{g_0}}$.}
\opt{B}{The time taken by the particle to move from end `A' to other `B' of the tunnel is $\dfrac{2\pi}{3}\sqrt{\dfrac{R}{g_0}}$.}
\opt{C}{The maximum speed of the particle inside the tunnel is $\sqrt{g_0R}$}
\opt{D}{The maximum speed of the particle inside the tunnel is $\sqrt{2g_0R}$}
\end{qoptions}
\end{bankq}
\begin{bankq}[id=EPS02-P1-PHY-S2-Q10, source={Enthusiast Pre-Score - Test 02 (11 Nov 2020)}, paper=1, qno=10, type=multi, subtopic={Satellites}, answer={B,D}, answer_src=key, difficulty=8, diff_parts={K2 S2 I2 C2 F1}, diff_note={concepts: bullet fired inward from a satellite, satellites | steps: energy and angular momentum; apogee; speed there; period}]
An earth satellite is revolving in a circular orbit of radius a with a velocity v$_0$. Mass of earth is M. A gun in the satellite is directly aimed toward earth. A bullet is fired from the gun with muzzle velocity $\dfrac{v_0}{2}$. Assume that mass of the satellite is very-very large with respect to the mass of the bullet. Choose the \textbf{\emph{CORRECT}} option(s) :
\begin{qoptions}
\opt{A}{Maximum distance of bullet from the centre of earth is $\dfrac{4a}{3}$.}
\opt{B}{Minimum speed of bullet (relative to earth) is $\dfrac{v_0}{2}$.}
\opt{C}{Minimum speed of bullet (relative to earth) is $\dfrac{3v_0}{4}$.}
\opt{D}{Time period of motion of bullet around earth is $T=\dfrac{2\pi}{\sqrt{GM}}\left(\dfrac{4a}{3}\right)^{3/2}$.}
\end{qoptions}
\end{bankq}
\begin{bankq}[id=EPS02-P1-PHY-S3-Q02, source={Enthusiast Pre-Score - Test 02 (11 Nov 2020)}, paper=1, qno=2, type=numerical, subtopic={Satellites}, answer={1.50}, answer_src=key, difficulty=7, diff_parts={K2 S2 I2 C1 F1}, diff_note={concepts: change of orbit by added energy, satellites | steps: initial energy; final energy; semi major axis relation}]
A particle of mass (m) is revolving around a planet of mass `M' in circular orbit of radius `3R'. It is given $\dfrac{GMm}{10aR}$ energy, so that now it moves in elliptical orbit of semi major axis `5R'. Find `a'.
\end{bankq}
\begin{bankq}[id=EPS02-P1-PHY-S3-Q05, source={Enthusiast Pre-Score - Test 02 (11 Nov 2020)}, paper=1, qno=5, type=numerical, subtopic={Kepler's laws}, answer={240.60}, answer_src=key, difficulty=8, diff_parts={K2 S2 I2 C2 F1}, diff_note={concepts: time between two points of an elliptical orbit, kepler's laws | steps: areal velocity; area swept; time ratio}]
Earth revolves around Sun with time period of `T = 365 days'. In the figure shown below, find the time taken (in days) by the Earth to move from A to B in the path ACB if the eccentricity of ellipse is 0.5.\par
\bfig{118.5}{EPS02_p1_phy_s3_q05_1.png}
\end{bankq}
\begin{bankq}[id=EPS02-P2-PHY-S1-Q01, source={Enthusiast Pre-Score - Test 02 (11 Nov 2020)}, paper=2, qno=1, type=multi, subtopic={Satellites}, answer={A,D}, answer_src=key, difficulty=8, diff_parts={K2 S2 I2 C2 F1}, diff_note={concepts: two satellites in nearby orbits, angular speeds | steps: angular speeds; difference; time between approaches}]
Two Earth's satellites move in a common plane along circular obrits. The orbital radius of one satellite is r while that of the other satellite is r $-$ $\Delta$r (Here $\Delta$r $\ll$ r).
\begin{qoptions}
\opt{A}{Time interval separates the periodic approaches of the satellites to each other over the minimum distance is $\dfrac{4\pi r^{5/2}}{3(GM)^{1/2}\Delta r}$.}
\opt{B}{Time interval separates the periodic approaches of the satellites to each other over the minimum distance is $\dfrac{2\pi r^{5/2}}{3(GM)^{1/2}\Delta r}$.}
\opt{C}{Angular velocity of approach between two satellites is $\dfrac{3(GM)^{1/2}\Delta r}{r^{5/2}}$.}
\opt{D}{Angular velocity of approach between two satellites is $\dfrac{3(GM)^{1/2}\Delta r}{2\,r^{5/2}}$.}
\end{qoptions}
\end{bankq}
\begin{bankq}[id=EPS02-P2-PHY-S3-Q01, source={Enthusiast Pre-Score - Test 02 (11 Nov 2020)}, paper=2, qno=1, type=numerical, subtopic={Kepler's laws}, answer={3.00}, answer_src=key, difficulty=7, diff_parts={K2 S2 I2 C1 F1}, diff_note={concepts: three body star system, kepler's laws | steps: force on a star; period relation; total mass}]
A trinary star system which is a system of three stars orbiting around centre of mass of system has time period of 3 years, while the distance between any two stars of the system is 2 astronomical unit. All the three stars are identical and mass of the sun is M and total mass of this multiple star system is $\dfrac ab\times M$. Then find $\dfrac b3$? (a and b are lowest possible integers)\par
(1 astronomical unit is equal to distance between sun and earth, time period of rotation of earth around sun is 1year)
\end{bankq}
\begin{bankq}[id=EPS02-P2-PHY-S3-Q04, source={Enthusiast Pre-Score - Test 02 (11 Nov 2020)}, paper=2, qno=4, type=numerical, subtopic={Escape velocity}, answer={2.00}, answer_src=key, difficulty=8, diff_parts={K2 S2 I2 C2 F1}, diff_note={concepts: escape speed from the midpoint of the earth and moon, escape velocity | steps: potential at the midpoint; energy condition; solve P}]
Given that mass of Earth is M and that of moon is m and distance between them is R. Neglect effect of any other planet or sun. A small mass is projected from midway from Earth and moon its escape velocity is $P\sqrt{\dfrac GR(m+M)}$ find P.
\end{bankq}

\section*{TOAS papers}
\begin{bankq}[id=TOAS08-P1-PHY-S1-Q05, source={TOAS - Test 08 (15 May 2022)}, paper=1, qno=5, type=multi, subtopic={Field and potential}, answer={C,D}, answer_src=key, difficulty=3, diff_parts={K1 S2 I0 C0 F1}, diff_note={concepts: gravitational field inside and outside a uniform sphere | steps: g inside is proportional to r and outside to 1/r\^{}2; potential inside and outside; pressure profile; pendulum period}]
Assuming the earth to be uniform sphere of radius $R$ and mass $M$ and neglecting rotational motion select correct option(s)
\begin{qoptions}
\opt{A}{\bfigi{103.8}{TOAS08_p1_phy_s1_q05_A.png}\newline $|\vec g|$ : Acceleration due to gravity, $r$ : distance from centre}
\opt{B}{\bfigi{94.5}{TOAS08_p1_phy_s1_q05_B.png}\newline $v$ : gravitational potential}
\opt{C}{\bfigi{79.8}{TOAS08_p1_phy_s1_q05_C.png}\newline $P$ : gravitational pressure.}
\opt{D}{\bfigi{83.3}{TOAS08_p1_phy_s1_q05_D.png}\newline $T$ : time period of simple pendulum.}
\end{qoptions}
\end{bankq}
\begin{bankq}[id=TOAS08-P1-PHY-S3-Q02, source={TOAS - Test 08 (15 May 2022)}, paper=1, qno=2, type=numerical, subtopic={Escape velocity}, answer={20.00}, answer_src=key, difficulty=5, diff_parts={K2 S2 I1 C0 F1}, diff_note={concepts: energy conservation, angular momentum conservation | steps: speed at the surface and critical impact parameter from energy and angular momentum; capture cross section; mass in the cloud volume | I: capture cross section from grazing orbits}]
A uniform cloud of dust particle approach a planet from a great distance with velocity $v_0$ relative to the planet. Radius of planet is $R$ and escape velocity from surface of the planet is $v_e=2v_0$. Length of the cloud along direction of motion is $\ell=R$, Cross-section of cloud perpendicular to it's direction of motion is very large and it's density is $\dfrac{4}{3000}$ times density of planet. If mass of planet $m$ is very large compared to the mass of cloud then mass of the dust particle that planet will collect during passage through the cloud is $\dfrac{m}{10x}$. The value of $x$ is :
\end{bankq}
\begin{bankq}[id=TOAS08-P2-PHY-S3-Q04, source={TOAS - Test 08 (15 May 2022)}, paper=2, qno=4, type=numerical, subtopic={Field and potential}, answer={29.53 to 29.55}, answer_src=key, difficulty=5, diff_parts={K2 S2 I1 C0 F1}, diff_note={concepts: angular momentum, energy conservation | steps: speeds at A and at the grazing point; conservation of angular momentum; conservation of energy; solve v | I: grazing means velocity tangential at R}]
A rocket is launched from point $A$, which is at distance $4R$ from centre of earth. If it passes grazing the earth surface then find value of $\dfrac{v}{100}$ where $v$ is velocity of projection. $\left[\text{use }\dfrac{GM}{R}=6.4\times10^7\,\mathrm{m^2/s^2}\right]$
\bfig{76.2}{TOAS08_p2_phy_s3_q04_1.png}
\end{bankq}
\begin{bankq}[id=TOAS08-P2-PHY-S3-Q07, source={TOAS - Test 08 (15 May 2022)}, paper=2, qno=7, type=numerical, subtopic={Satellites}, answer={58.00}, answer_src=key, difficulty=5, diff_parts={K2 S2 I1 C0 F1}, diff_note={concepts: collision with a satellite, new orbit | steps: momentum conservation for the stick; speed of the satellite and meteorite combination; energy and angular momentum for the new orbit; solve for n | I: new orbit conditions at the periapsis}]
A meteorite is approaching a planet in a straight line passing through the centre of planet. It collides with a satellite orbiting the planet in orbit of radius $r$ having mass ten times of mass of meteorite and sticks with it. As a result of collision satellite goes to new orbit with minimum distance $\dfrac r2$ from the centre of planet. If speed of meteorite just before collision is $u=\sqrt{\dfrac{nGM}{R}}$. Find value of $n$.
\end{bankq}
\begin{bankq}[id=TOAS08-P2-PHY-S3-Q08, source={TOAS - Test 08 (15 May 2022)}, paper=2, qno=8, type=numerical, subtopic={Kepler's laws}, answer={2.25}, answer_src=key, difficulty=3, diff_parts={K1 S1 I1 C0 F1}, diff_note={concepts: binary star motion | steps: both stars share the angular velocity; areal speed ratio from the squares of the radii | I: radius ratio inversely proportional to mass}]
Figure shows a binary star system revolving about their COM. The masses of star $A$ \& $B$ are $15\times10^{30}$ kg and $30\times10^{30}$ kg respectively. Find the ratio of area swept by star $A$ to area swept by star $B$ in a common time interval.
\bfig{99.7}{TOAS08_p2_phy_s3_q08_1.png}
\end{bankq}
\begin{bankq}[id=TOAS09-P1-PHY-S2-Q08, source={TOAS - Test 09 (30 May 2022)}, paper=1, qno=8, type=multi, subtopic={Satellites}, answer={A,C}, answer_src=key, difficulty=4, diff_parts={K2 S2 I0 C0 F1}, diff_note={concepts: relative angular speed of a satellite, Kepler's third law | steps: relative period of the satellite and the earth; geostationary orbit height; orbit radius scaling}, also_in={TNAS - Test 02 (07 May 2023) P1 Q9}]
Consider a satellite `S' which is just above the ECUADOR and rotating on equatorial plane of earth. The sense of rotation of earth and satellite are similar means both are rotating from west to east. The time period of satellite is 6 hours. At $t=0$, satellite is just above the Brazil. Now choose the correct option (s).
\bfig{130.6}{TOAS09_p1_phy_s2_q08_1.png}
\begin{qoptions}
\opt{A}{The time after which the satellite will be above the ECUADOR again is 8 hours.}
\opt{B}{The above satellite may be a geostationary satellite.}
\opt{C}{The height of geostationary satellite above the sea level is 35,800 km approximately.}
\opt{D}{If the distance from the centre of earth to the satellite `S' becomes $4r_0$, then the time after which the satellite will be above the ECUADOR again is 24 hours.}
\end{qoptions}
\end{bankq}
\begin{bankq}[id=TOAS09-P1-PHY-S2-Q09, source={TOAS - Test 09 (30 May 2022)}, paper=1, qno=9, type=multi, subtopic={Field and potential}, answer={B,C}, answer_src=key, difficulty=5, diff_parts={K2 S2 I1 C0 F1}, diff_note={concepts: motion in a tunnel along a chord, SHM | steps: gravity component along the chord is -g x/R; omega = sqrt(g/R); amplitude from the initial speed; time between ends and maximum speed | I: chord motion is SHM}]
A tunnel is dug across the earth along a chord AB as shown, A particle is projected inside the tunnel from end `A' with an initial velocity $v_0=\dfrac{\sqrt{3g_0R}}{2}$, there is no friction between the tunnel and the particle. Then choose the correct option (s), (where $g_0$ = gravitational acceleration on the surface of Earth, $R$ = radius of earth)
\bfig{124.6}{TOAS09_p1_phy_s2_q09_1.png}
\begin{qoptions}
\opt{A}{The time taken by the particle to move from end `A' to other end `B' of the tunnel is $\dfrac{\pi}{3}\sqrt{\dfrac{R}{g_0}}$.}
\opt{B}{The time taken by the particle to move from end `A' to other end `B' of the tunnel is $\dfrac{2\pi}{3}\sqrt{\dfrac{R}{g_0}}$.}
\opt{C}{The maximum speed of the particle inside the tunnel is $\sqrt{g_0R}$.}
\opt{D}{The maximum speed of the particle inside the tunnel is $\sqrt{2g_0R}$.}
\end{qoptions}
\end{bankq}
\begin{bankq}[id=TOAS09-P1-PHY-S4-Q07, source={TOAS - Test 09 (30 May 2022)}, paper=1, qno=7, type=numerical, subtopic={Escape velocity}, answer={7}, answer_src=key, difficulty=3, diff_parts={K1 S1 I1 C0 F1}, diff_note={concepts: time of flight at escape speed | steps: v = sqrt(2GM/r) at escape; t = integral of dr/v from R to 4R | I: escape speed gives total energy zero}]
A small body of mass `m' is projected from the surface of a given planet as shown in the figure. If small body has velocity which is just sufficient to get out from the gravity of planet then time taken by small body to reach at height of $3R$ from the surface of planet is $\dfrac{K}{3}\sqrt{\dfrac{2R^3}{GM}}$. Find the value of K.
\bfig{79.1}{TOAS09_p1_phy_s4_q07_1.png}
\end{bankq}
\begin{bankq}[id=TOAS09-P2-PHY-S1-Q07, source={TOAS - Test 09 (30 May 2022)}, paper=2, qno=7, type=single, subtopic={Field and potential}, answer={D}, answer_src=key, difficulty=2, diff_parts={K1 S1 I1 C0 F0}, diff_note={concepts: gravitational field with cavities | steps: superposition of the sphere field and negative mass cavities; symmetric cancellation | I: cavities are symmetric}]
A solid sphere of uniform density \& having total mass $(M)$ now two spherical cavity is removed as shown in figure. Radius of sphere is 4 meter. Radius of each spherical cavity is 1 meter. Value of gravitation field at centre of sphere will be :-
\bfig{127.3}{TOAS09_p2_phy_s1_q07_1.png}
\begin{qoptions}
\opt{A}{$\dfrac{31GM}{1024}$}
\opt{B}{$\dfrac{31GM}{256}$}
\opt{C}{$\dfrac{GM}{1024}$}
\opt{D}{Zero}
\end{qoptions}
\end{bankq}
\begin{bankq}[id=TOAS09-P2-PHY-S1-Q08, source={TOAS - Test 09 (30 May 2022)}, paper=2, qno=8, type=single, subtopic={Field and potential}, answer={A}, answer_src=key, difficulty=1, diff_parts={K1 S1 I0 C0 F0}, diff_note={concepts: orbit of four equal masses | steps: net gravitational force on one mass; centripetal equation}]
Four similar masses orbiting around a circle of radius `r'. Due to their mutual gravitation attraction force velocity of each particle will be :-
\bfig{92.0}{TOAS09_p2_phy_s1_q08_1.png}
\begin{qoptions}
\opt{A}{$\sqrt{\dfrac{Gm}{r}\left(\dfrac{1+2\sqrt2}{4}\right)}$}
\opt{B}{$\sqrt{\dfrac{Gm}{r}\left(1+2\sqrt2\right)}$}
\opt{C}{$\sqrt{\dfrac{Gm}{r}}$}
\opt{D}{$\sqrt{\dfrac{Gm}{r}\left(\dfrac{1+2\sqrt2}{2}\right)}$}
\end{qoptions}
\end{bankq}
\begin{bankq}[id=TOAS09-P2-PHY-S1-Q19, source={TOAS - Test 09 (30 May 2022)}, paper=2, qno=19, type=single, subtopic={Escape velocity}, answer={A}, answer_src=key, difficulty=3, diff_parts={K2 S2 I0 C0 F0}, diff_note={concepts: potential inside a uniform sphere, escape speed | steps: potential at depth R/4 inside the sphere; energy conservation; escape speed}]
A particle is projected from bottom of a blind narrow tunnel of depth $\dfrac R4$ made along a radius inside the earth. If minimum velocity needed to escape out of the gravitational field of the earth is $\dfrac14\sqrt{\dfrac{xGM}{R}}$. Find $x$. $M$ and $R$ are mass and radius of the earth respectively.
\begin{qoptions}
\opt{A}{39}
\opt{B}{78}
\opt{C}{86}
\opt{D}{0}
\end{qoptions}
\end{bankq}
\begin{bankq}[id=TOAS09-P2-PHY-S2-Q01, source={TOAS - Test 09 (30 May 2022)}, paper=2, qno=1, type=numerical, subtopic={Field and potential}, answer={11}, answer_src=solved, working={Energy gives $v_{top}^2=0.1gR$; angular momentum $v\sin\theta R=v_{top}\cdot\frac54R$ gives $\cos\theta=\sqrt{11}/4$}, difficulty=5, diff_parts={K2 S2 I1 C0 F1}, diff_note={concepts: energy, angular momentum for a projectile in gravity | steps: speed at the top from energy; angular momentum conservation; angle of projection | I: apogee condition}]
A body is thrown from surface of earth with velocity $\left(\sqrt{\dfrac{gR_e}{2}}\right)$ [$R_e$ radius of earth] at some angle from vertical. Maximum height reach by the body is $\dfrac{R_e}{4}$. If angle of projection is $\cos^{-1}\!\left[\dfrac{\sqrt x}{4}\right]$ find the value of x.
\end{bankq}
\begin{bankq}[id=TOAS09-P2-PHY-S2-Q02, source={TOAS - Test 09 (30 May 2022)}, paper=2, qno=2, type=numerical, subtopic={Field and potential}, answer={2}, answer_src=solved, working={Gravitational pressure $P(r)=\frac{2\pi G\rho^2}3(R^2-r^2)$: ratio at R/2 is 3/4, so $3x/8=3/4$, $x=2$}, difficulty=3, diff_parts={K1 S2 I0 C0 F1}, diff_note={concepts: gravitational pressure in a uniform sphere | steps: g(r) is proportional to r; integrate rho g dr; ratio at R/2}]
A uniform solid sphere having mass density is $\rho$. The ratio of gravitational pressure at a distance $r=R/2$ to the centre of sphere is $\dfrac{3x}{8}$, the value of x will be
\bfig{60.7}{TOAS09_p2_phy_s2_q02_1.png}
\end{bankq}
\begin{bankpassage}[id=TOAS10-P1-PHY-S2-G11-12, source={TOAS - Test 10 (12 Jun 2022)}, questions={11-12}]
A rocket is launched from and returns to a spherical planet of radius $R$ in such a way that its velocity vector on return is parallel to its launch vector. The angular separation at the centre of the planet between the launch and arrival points is $\theta$.
\end{bankpassage}
\begin{bankq}[id=TOAS10-P1-PHY-S2-Q11, source={TOAS - Test 10 (12 Jun 2022)}, paper=1, qno=11, type=multi, subtopic={Satellites}, answer={A,B,C}, answer_src=key, difficulty=5, diff_parts={K2 S2 I1 C0 F1}, diff_note={concepts: conic orbits, symmetry of the trajectory | steps: conic orbit with the planet centre at a focus; parallel velocities imply symmetry; the maximum height consistent with the condition | I: parallel launch and arrival velocity vectors}, group=TOAS10-P1-PHY-S2-G11-12]
Which of the following can be the distance of the rocket above the surface of the planet ?
\begin{qoptions}
\opt{A}{$R$}
\opt{B}{$\dfrac R2$}
\opt{C}{$\dfrac R3$}
\opt{D}{$4R$}
\end{qoptions}
\end{bankq}
\begin{bankq}[id=TOAS10-P1-PHY-S2-Q12, source={TOAS - Test 10 (12 Jun 2022)}, paper=1, qno=12, type=multi, subtopic={Satellites}, answer={A,B}, answer_src=key, difficulty=6, diff_parts={K2 S2 I2 C0 F1}, diff_note={concepts: orbit geometry, Kepler's equation | steps: semi-major axis from the launch condition; Kepler's equation for the time; evaluate for theta -> 0 and theta = 90 degrees | I: orbit symmetric about the apogee}, group=TOAS10-P1-PHY-S2-G11-12]
How long does the flight of rocket take, if the period of satellite flying around the planet just above its surface is $T_0$ ?
\begin{qoptions}
\opt{A}{$T_0\left(\dfrac12+\dfrac1\pi\right)$ if $\theta\to0$}
\opt{B}{$T_0\left(\dfrac12+\dfrac{1}{\sqrt2\pi}\right)$ if $\theta=90^\circ$}
\opt{C}{$T_0\left(\dfrac12+\dfrac2\pi\right)$ if $\theta\to0$}
\opt{D}{$T_0\left(\dfrac12+\dfrac{\sqrt2}{\pi}\right)$ if $\theta=90^\circ$}
\end{qoptions}
\end{bankq}
\begin{bankq}[id=TOAS10-P1-PHY-S3-Q05, source={TOAS - Test 10 (12 Jun 2022)}, paper=1, qno=5, type=numerical, subtopic={Escape velocity}, answer={3.00}, answer_src=key, difficulty=4, diff_parts={K2 S2 I0 C0 F1}, diff_note={concepts: energy, angular momentum in a planetary field | steps: escape speed and launch speed at 2R; energy conservation; angular momentum conservation at the farthest point; solve for the radius}]
A particle of mass $m$ is fired from the point A (as shown in the figure) situated at a distance of $2R$ from the centre of the planet with a velocity $\dfrac{1}{\sqrt2}$ times the escape velocity from point A. The radius of the planet is $R$. The maximum distance of the particle from the centre in the subsequent motion is $nR$. Find the value of $n$.
\bfig{62.0}{TOAS10_p1_phy_s3_q05_1.png}
\end{bankq}
\begin{bankq}[id=TOAS11-P1-PHY-S1-Q04, source={TOAS - Test 11 (27 Jun 2022)}, paper=1, qno=4, type=single, subtopic={Field and potential}, answer={B}, answer_src=key, difficulty=2, diff_parts={K1 S1 I1 C0 F0}, diff_note={concepts: projectile in a radial gravitational field | steps: angular momentum conservation about the earth centre; apex speed from the launch speed | I: speed at the apex from angular momentum}]
A projectile is thrown at angle $60^\circ$ with surface of earth, with velocity $v_0$ such that it attains a maximum height equal to radius of earth. Then :-
\bfig{66.1}{TOAS11_p1_phy_s1_q04_1.png}
\begin{qoptions}
\opt{A}{minimum velocity would be $\dfrac{V_0}{2}$}
\opt{B}{minimum velocity would be $\dfrac{V_0}{4}$}
\opt{C}{minimum velocity would be $\dfrac{V_0}{3}$}
\opt{D}{minimum velocity would be $\dfrac{V_0}{6}$}
\end{qoptions}
\end{bankq}
\begin{bankq}[id=TOAS11-P2-PHY-S1-Q07, source={TOAS - Test 11 (27 Jun 2022)}, paper=2, qno=7, type=multi, subtopic={Satellites}, answer={C}, answer_src=key, difficulty=2, diff_parts={K1 S1 I0 C0 F1}, diff_note={concepts: satellites in the same orbit | steps: speed independent of mass; potential energy and kinetic energy scale with mass}]
Two satellites A and B move round the earth in the same orbit. The mass of B is three times the mass of A.
\begin{qoptions}
\opt{A}{Speeds of A and B are different}
\opt{B}{The potential energy of earth $+$A is same as that of earth $+$B.}
\opt{C}{The kinetic energy of A and B are different.}
\opt{D}{The total energy of earth $+$A is same as that of earth $+$B.}
\end{qoptions}
\end{bankq}
\begin{bankq}[id=TOAS11-P2-PHY-S2-Q11, source={TOAS - Test 11 (27 Jun 2022)}, paper=2, qno=11, type=single, subtopic={Satellites}, answer={B}, answer_src=key, difficulty=4, diff_parts={K2 S2 I1 C0 F0}, diff_note={concepts: orbital mechanics, angular momentum conservation | steps: orbit conditions after each change; energy of the satellite; coverage; angular momentum of the planet | I: angular momentum of the planet is conserved when it contracts}]
\textbf{Answer the following by appropriately matching the lists based on the information given in the paragraph.}\par
A planet is rotating about it's axis. The radius of planet is $R$ and its mass is $m$. It has a geostationary satellite in equatorial plane at a height $h$ from surface.\par
\smallskip
\noindent{\footnotesize\setlength{\tabcolsep}{2pt}\begin{tabular}{|L{.95cm}|L{2.7cm}|L{.95cm}|L{3.1cm}|}\hline
\multicolumn{2}{|c|}{List-I} & \multicolumn{2}{c|}{List-II}\\\hline
(I) & If the planet suddenly contracts without change in it's mass & (P) & The energy of satellite decreases.\\\hline
(II) & If the planet suddenly stop rotating. & (Q) & The energy of satellite increases.\\\hline
(III) & If the geostationary satellite is pushed in direction of motion to suddenly double the velocity. & (R) & The energy of satellite remains same.\\\hline
(IV) & If the geostationary satellite suddenly detaches half of it's mass without any relative velocity. & (S) & The satellite covers larger area on the planet.\\\hline
 & & (T) & The satellite does not remain geostationary.\\\hline
 & & (U) & The satellite should have now been kept at a farther orbit from earth to be geostationary.\\\hline
\end{tabular}}
\begin{qoptions}
\opt{A}{I $\to$ R,S,T;II $\to$ R,S,T;III $\to$ Q,S,T;IV $\to$ P}
\opt{B}{I $\to$ R,S,T;II $\to$ R,S,T;III $\to$ Q,S,T;IV $\to$ Q}
\opt{C}{I $\to$ R,T,U;II $\to$ R,S,T;III $\to$ P,R,S,T;IV $\to$ R}
\opt{D}{I $\to$ R,S,T;II $\to$ R,S;III $\to$ Q,S,T;IV $\to$ P}
\end{qoptions}
\end{bankq}
\begin{bankq}[id=TOAS11-P2-PHY-S2-Q12, source={TOAS - Test 11 (27 Jun 2022)}, paper=2, qno=12, type=single, subtopic={Satellites}, answer={A}, answer_src=key, difficulty=2, diff_parts={K2 S1 I0 C0 F0}, diff_note={concepts: orbital mechanics, circular orbit relations | steps: angular velocity, energy, angular momentum and speed in terms of G, M, R + h; dependencies}]
\textbf{Answer the following by appropriately matching the lists based on the information given in the paragraph.}\par
A planet is rotating about it's axis. The radius of planet is $R$ and its mass is $m$. It has a geostationary satellite in equatorial plane at a height $h$ from surface.\par
\smallskip
\noindent{\footnotesize\setlength{\tabcolsep}{2pt}\begin{tabular}{|L{.95cm}|L{2.7cm}|L{.95cm}|L{3.1cm}|}\hline
\multicolumn{2}{|c|}{List-I} & \multicolumn{2}{c|}{List-II}\\\hline
(I) & Angular velocity of satellite. & (P) & Proportional to $\sqrt{R+h}$\\\hline
(II) & Energy of satellite & (Q) & Proportional to $\dfrac{1}{\sqrt{R+h}}$\\\hline
(III) & Angular momentum of satellite about centre of planet & (R) & Depends on mass of satellite.\\\hline
(IV) & Speed of satellite & (S) & Depends on mass of planet.\\\hline
 & & (T) & Proportional to $\sqrt G$\\\hline
 & & (U) & Depends on angular velocity of planet.\\\hline
\end{tabular}}
\begin{qoptions}
\opt{A}{I $\to$ S,T,U;II $\to$ R,S,U;III $\to$ P,R,S,T,U;IV $\to$ Q,S,T,U}
\opt{B}{I $\to$ S,T,U;II $\to$ R,S;III $\to$ P,R,S,T;IV $\to$ Q,S,T}
\opt{C}{I $\to$ R,T,U;II $\to$ R,S;III $\to$ P,R,S,T;IV $\to$ Q,T,U}
\opt{D}{I $\to$ P,S,T;II $\to$ R,S,U;III $\to$ P,R,S,T,U;IV $\to$ Q,R,S,T}
\end{qoptions}
\end{bankq}
\begin{bankq}[id=TOAS11-P2-PHY-S3-Q01, source={TOAS - Test 11 (27 Jun 2022)}, paper=2, qno=1, type=numerical, subtopic={Escape velocity}, answer={2.00}, answer_src=key, difficulty=4, diff_parts={K1 S2 I1 C0 F1}, diff_note={concepts: escape from an elliptic path | steps: energy at the launch point; speed at infinity; ratio to v0 | I: escape speed from the radius of the circular orbit}]
An Earth satellite revolves on a circular orbit at a linear speed $v_0$. A missile device increases the absolute value of this speed to $1.5v_0$, and the satellite goes into an unknown path as shown (Fig). With what speed will the satellite leave the gravitational field of earth? The resistance of the cosmic dust is not taken into account. Mass of earth is $M$. If your answer is $\dfrac{v_0}{\alpha}$, fill $\alpha$ in OMR sheet.
\bfig{111.5}{TOAS11_p2_phy_s3_q01_1.png}
\end{bankq}
\begin{bankq}[id=TOAS12-P1-PHY-S1-Q07, source={TOAS - Test 12 (10 Jul 2022)}, paper=1, qno=7, type=multi, subtopic={Satellites}, answer={A,C}, answer_src=key, difficulty=6, diff_parts={K2 S2 I2 C0 F1}, diff_note={concepts: dimensional analysis, orbit decay by radiation | steps: dimensions fix the exponents of c and G; orbital energy; rate of change of period from dE/dt | I: power law from dimensions}]
A massive star of mass $M$ is in uniform circular orbit around a supermassive black hole of mass $M_b$. Initially, the radius and angular frequency of the orbit are $R$ and $\omega$ respectively. According to Einstein's theory of general relativity the space around the two objects is distorted and gravitational waves are radiated. Energy is lost through this radiation and as a result the orbit of the star shrinks gradually. One may assume, however, that the orbit remains circular throughout and Newtonian mechanics holds. Choose the correct option(s) :-
\begin{qoptions}
\opt{A}{The power radiated through gravitational wave by this star is given by $L_G=KC^xG^yM^2R^4\omega^6$ where $c$ is the speed of light, $G$ is the universal gravitational constant, and $K$ is a dimensionless constant, then $x=-5$ \& $y=1$}
\opt{B}{Total mechanical energy of the star-black hole system is $\dfrac{GMm_b}{2R}$}
\opt{C}{An expression for the rate of decrease in the orbital period ($dT/dt$) in terms of the masses, period $T$ and constants given by $\dfrac{dT}{dt}=-\dfrac{3KG^{5/3}M_b^{2/3}M(2\pi)^{8/3}}{c^5T^{5/3}}$}
\opt{D}{An expression for the rate of decrease in the orbital period ($dT/dt$) in terms of the masses, period $T$ and constants given by $\dfrac{dT}{dt}=\dfrac{3KG^{5/3}M_b^{2/3}M}{c^5T^{5/3}}$}
\end{qoptions}
\end{bankq}
\begin{bankq}[id=TOAS12-P2-PHY-S3-Q06, source={TOAS - Test 12 (10 Jul 2022)}, paper=2, qno=6, type=numerical, subtopic={Field and potential}, answer={30.76}, answer_src=key, difficulty=4, diff_parts={K1 S2 I1 C0 F1}, diff_note={concepts: two bodies in mutual attraction | steps: momentum and energy conservation; speed of B gives speed of A; energy equation for the separation | I: centre of mass fixed}]
Two particles A and B of mass 1kg and 2kg respectively are kept 1 meter apart and are released to move under mutual attraction. When speed of B is 3.6 cm/hour, what is the separation (in cm) between the particles. Take $G=\dfrac{20}{3}\times10^{-11}$ Nm$^2$/kg.
\end{bankq}
\begin{bankq}[id=TOAS15-P1-PHY-S1-Q06, source={TOAS - Test 15 (22 Aug 2022)}, paper=1, qno=6, type=single, subtopic={Field and potential}, answer={A}, answer_src=key, difficulty=4, diff_parts={K2 S2 I1 C0 F0}, diff_note={concepts: energy conservation in a self-gravitating shell, field of a shell | steps: a surface particle feels half of GM/R\^{}2; energy per unit mass; turning point where v = 0}]
An asteroid of mass $M$ exploded into a spherical homogeneous cloud [spherical shell] in free space. Due to energy received by the explosion, the cloud expands and the expansion is spherically symmetric. At an instant, when radius of the cloud is $R_0$, all of its particles on the surface are observed receding radially away from the centre of the cloud with a velocity $v_0$. What will be the radius of the cloud, when its expansion ceases?
\begin{qoptions}
\opt{A}{$\dfrac{GMR_0}{GM-v_0^2R_0}$}
\opt{B}{$\dfrac{GM}{v_0^2}$}
\opt{C}{$\dfrac{2GMR_0}{2GM-v_0^2R_0}$}
\opt{D}{$\dfrac{GMR_0}{2GM-V_0^2R_0}$}
\end{qoptions}
\end{bankq}
\begin{bankq}[id=TOAS15-P1-PHY-S3-Q03, source={TOAS - Test 15 (22 Aug 2022)}, paper=1, qno=3, type=numerical, subtopic={Field and potential}, answer={80.00}, answer_src=key, difficulty=6, diff_parts={K2 S2 I1 C1 F1}, diff_note={concepts: angular momentum conservation, energy conservation in gravity | steps: tangential speed at the surface from L; energy equation; solve for the initial speed; x}]
A particle is projected from point, that is at a distance $4R$ from centre of the earth, with speed $50\sqrt2\,x$ m/s, in a direction making angle $30^\circ$ with vertical, away from the earth. Find $x$ if particle passes grazing the surface of the earth use $\dfrac{GM}{R}=6.4\times10^7$m$^2$/s$^2$, where $R$ is radius of earth and $M$ is mass of the earth.
\end{bankq}
\begin{bankq}[id=TOAS15-P2-PHY-S1-Q03, source={TOAS - Test 15 (22 Aug 2022)}, paper=2, qno=3, type=multi, subtopic={Field and potential}, answer={A,B,D}, answer_src=solved, working={P lies between the shells: the outer shell gives no field there, the inner shell attracts $m_3$ toward O with $Gm_1m_3/r^2$; $m_2$ plays no role}, difficulty=3, diff_parts={K1 S2 I0 C0 F1}, diff_note={concepts: field of concentric shells | steps: inside the outer shell its field vanishes; inner shell acts as a point mass; test each statement}]
Two concentric spherical shells masses $m_1$ and $m_2$ and radii $r_1$ and $r_2$. Then
\bfig{82.5}{TOAS15_p2_phy_s1_q03_1.png}
\begin{qoptions}
\opt{A}{outer shell will have no contribution in gravitational field at point P.}
\opt{B}{force on P is directed towards O.}
\opt{C}{force on P is $\dfrac{Gm_1m_2}{r^2}$.}
\opt{D}{force on P is $\dfrac{Gm_1m_3}{r^2}$.}
\end{qoptions}
\end{bankq}
\begin{bankq}[id=TOAS15-P2-PHY-S3-Q05, source={TOAS - Test 15 (22 Aug 2022)}, paper=2, qno=5, type=numerical, subtopic={Escape velocity}, answer={2}, answer_src=key, difficulty=5, diff_parts={K2 S2 I1 C0 F1}, diff_note={concepts: escape energy, energy conservation in gravity | steps: kinetic energy is one third of the escape value; energy equation for the height; read N}]
A particle is projected vertically upwards from the surface of the earth with a kinetic energy equal to $\dfrac13$ times the minimum kinetic energy needed to escape. If radius of the earth is 6400km. The maximum height attained by the particle from the surface of the earth is $\dfrac1N$ times the radius of earth. Determine the value on $N$.
\end{bankq}
\begin{bankq}[id=TOAS17-P1-PHY-S3-Q03, source={TOAS - Test 17 (18 Sep 2022)}, paper=1, qno=3, type=numerical, subtopic={Escape velocity}, answer={8.00}, answer_src=key, difficulty=6, diff_parts={K2 S2 I1 C1 F1}, diff_note={concepts: jump height and escape speed, gravity of a planet | steps: height gives launch speed; escape speed in terms of radius and density; compare with Earth; radius}]
A person can throw a ball to a maximum height of 20 m on the earth. If same person throw same ball on an imaginary planet whose density is twice of the earth then maximum possible radius of planet so that ball can escape from planet is $R$ km. Find $R$. (Radius of the Earth = $6.4\times10^6$ m)
\end{bankq}
\begin{bankq}[id=TOAS17-P2-PHY-S2-Q02, source={TOAS - Test 17 (18 Sep 2022)}, paper=2, qno=2, type=numerical, subtopic={Field and potential}, answer={7}, answer_src=key, difficulty=5, diff_parts={K1 S2 I1 C1 F1}, diff_note={concepts: gravitational field inside a layered planet | steps: field at the surface; field at depth R; equate; density ratio}]
A planet is made of two materials of density `$\rho_1$' and `$\rho_2$' as shown in figure. Acceleration due to gravity at surface of planet is same as that at depth `$R$'. The ratio $\dfrac{\rho_1}{\rho_2}$ is equal to $\dfrac x3$. Find the value of `$x$'.
\bfig{128.3}{TOAS17_p2_phy_s2_q02_1.png}
\end{bankq}
\begin{bankq}[id=TOAS18-P1-PHY-S3-Q08, source={TOAS - Test 18 (30 Oct 2022)}, paper=1, qno=8, type=numerical, subtopic={Satellites}, answer={16}, answer_src=solved, working={Kepler third law with total mass $4M$ and separation 4 AU: $4M=4^3/1^2=64$ solar masses, $M=16\,M_\odot$}, difficulty=2, diff_parts={K1 S1 I0 C0 F1}, diff_note={concepts: binary star orbit | steps: Kepler law for the total mass; ratio}]
Two starts are revolving in circular paths about center of mass of system. If there masses are $M$ and $3M$ and separation between them is 4 times separation between the Earth and Sun. If the orbital period is 1 year find ratio of $M$ and mass of Sun.
\end{bankq}
\begin{bankq}[id=TOAS18-P2-PHY-S1-Q02, source={TOAS - Test 18 (30 Oct 2022)}, paper=2, qno=2, type=single, subtopic={Field and potential}, answer={D}, answer_src=key, difficulty=5, diff_parts={K2 S2 I1 C1 F0}, diff_note={concepts: motion in a chord tunnel, normal force | steps: gravity inside the earth; component along the tunnel; normal reaction; plot}]
A smooth tunnel is made along chord of earth at $\dfrac R2$ from centre of earth ($R$ = radius of earth). A object is released from one end of tunnel. Which of the following is correct graph of variation of contact force $N$ vs distance from centre of tunnel :
\begin{qoptions}
\opt{A}{\bfigi{75.7}{TOAS18_p2_phy_s1_q02_A.png}}
\opt{B}{\bfigi{75.2}{TOAS18_p2_phy_s1_q02_B.png}}
\opt{C}{\bfigi{70.5}{TOAS18_p2_phy_s1_q02_C.png}}
\opt{D}{\bfigi{75.2}{TOAS18_p2_phy_s1_q02_D.png}}
\end{qoptions}
\end{bankq}

\section*{TAAS papers}
\begin{bankq}[id=TAAS01-P1-PHY-S3-Q01, source={TAAS - Test 01 (18 Jun 2023)}, paper=1, qno=1, type=numerical, subtopic={Satellites}, answer={3.00}, answer_src=key, difficulty=6, diff_parts={K2 S2 I1 C1 F1}, diff_note={concepts: tidal locking, weightlessness | steps: condition for zero weight on the facing point; orbital angular velocity; expansion for small r; x}]
A small spherical uniform moon of radius $r$ and mass $m$ is orbiting circularly around a planet (uniform sphere) of mass $M$ while keeping the same face towards the planet. Any object on the moon's surface facing the planet experiences weightlessness if radius of the orbit of moon is $r\sqrt[3]{\dfrac{xM}m}$ where $x$ is \underline{\hspace{1.5cm}}
\end{bankq}
\begin{bankq}[id=TAAS02-P1-PHY-S3-Q08, source={TAAS - Test 02 (30 Jul 2023)}, paper=1, qno=8, type=single, subtopic={Satellites}, answer={C}, answer_src=key, difficulty=2, diff_parts={K1 S2 I0 C0 F0}, diff_note={concepts: orbit type with the launch speed | steps: circular speed; escape speed; classify each speed range}]
A particle is taken to a distance $r$ ($>R$) from centre of earth. $R$ is the radius of the earth. It is given velocity $v$ which is perpendicular to radius. With the given values of $v$ in List-I, you have to match the resultant path of particle in List-II.
\begin{center}\footnotesize\setlength{\tabcolsep}{2pt}
\begin{tabular}{|L{0.6cm}L{2.5cm}|L{0.6cm}L{1.8cm}|}\hline
\multicolumn{2}{|c|}{List-I}&\multicolumn{2}{c|}{List-II}\\\hline
(P)&$v=\sqrt{\dfrac{GM}r}$&(1)&Elliptical\\\hline
(Q)&$v=\sqrt{\dfrac{2GM}r}$&(2)&Parabolic\\\hline
(R)&$v>\sqrt{\dfrac{2GM}r}$&(3)&Hyperbolic\\\hline
(S)&$\sqrt{\dfrac{GM}r}<v<\sqrt{\dfrac{2GM}r}$&(4)&Circular\\\hline
&&(5)&Straight line\\\hline
\end{tabular}\end{center}
\begin{qoptions}
\opt{A}{P $\to$ 3;Q $\to$ 1;R $\to$ 4;S $\to$ 2}
\opt{B}{P $\to$ 2;Q $\to$ 4;R $\to$ 3;S $\to$ 1}
\opt{C}{P $\to$ 4;Q $\to$ 2;R $\to$ 3;S $\to$ 1}
\opt{D}{P $\to$ 1;Q $\to$ 2;R $\to$ 3;S $\to$ 4}
\end{qoptions}
\end{bankq}

\section*{TMAS papers}
\begin{bankq}[id=TMAS01-P1-PHY-S1-Q13, source={TMAS - Test 01 (09 Jul 2023)}, paper=1, qno=13, type=single, subtopic={Kepler's laws}, answer={C}, answer_src=key, difficulty=2, diff_parts={K1 S1 I0 C1 F0}, diff_note={concepts: Kepler third law | steps: semi major axis; relation to period}]
The earth is moving on an elliptical path, whose one focus sun is situated as shown in figure. If $AS=r_{\min}$ and $SB=r_{\max}$, the sun and earth system obey the Kepler's law, the square of time-period is directly proportional to
\bfig{90.3}{TMAS01_p1_phy_s1_q13_1.png}
\begin{qoptions}
\opt{A}{$(r_{\min})^3$}
\opt{B}{$(r_{\max})^3$}
\opt{C}{$\left(\dfrac{r_{\min}+r_{\max}}2\right)^3$}
\opt{D}{$\left(\dfrac{r_{\min}r_{\max}}{r_{\min}+r_{\max}}\right)^3$}
\end{qoptions}
\end{bankq}
\begin{bankq}[id=TMAS01-P1-PHY-S1-Q14, source={TMAS - Test 01 (09 Jul 2023)}, paper=1, qno=14, type=single, subtopic={Satellites}, answer={A}, answer_src=key, difficulty=5, diff_parts={K2 S2 I1 C1 F0}, diff_note={concepts: satellite with drag, energy | steps: orbit decay; speed increases; angular momentum falls; period}]
When a satellite in a circular orbit around the earth enters the atmospheric region, it encounters small air resistance to its motion. Then
\begin{qoptions}
\opt{A}{It kinetic energy increases}
\opt{B}{It kinetic energy decreases}
\opt{C}{Its angular momentum about the earth increases}
\opt{D}{Its period of revolution around the earth increases}
\end{qoptions}
\end{bankq}
\begin{bankq}[id=TMAS01-P1-PHY-S1-Q15, source={TMAS - Test 01 (09 Jul 2023)}, paper=1, qno=15, type=single, subtopic={Satellites}, answer={B}, answer_src=key, difficulty=4, diff_parts={K1 S2 I1 C1 F0}, diff_note={concepts: binary stars | steps: orbital radius; gravitational force; period}]
Two spherical stars each of mass $m$ with their centres at distance $D$ revolve under mutual gravitational attraction about the point midway between their centres. The period of revolution will be :-
\begin{qoptions}
\opt{A}{$\pi\sqrt{\dfrac{D^3}{GM}}$}
\opt{B}{$2\pi\sqrt{\dfrac{D^3}{2GM}}$}
\opt{C}{$2\pi\sqrt{\dfrac{2D^3}{GM}}$}
\opt{D}{$\pi\sqrt{\dfrac{GM}{D^3}}$}
\end{qoptions}
\end{bankq}
\begin{bankq}[id=TMAS01-P1-PHY-S1-Q16, source={TMAS - Test 01 (09 Jul 2023)}, paper=1, qno=16, type=single, subtopic={Satellites}, answer={A}, answer_src=key, difficulty=3, diff_parts={K1 S1 I1 C1 F0}, diff_note={concepts: orbital period relation | steps: T squared proportional to r cubed; differentiate}]
A satellite is revolving round the earth in an orbit of radius $r$ with time period $T$. If the satellite is revolving round the earth in an orbit of radius $r+\Delta r$ ($\Delta r\ll r$) with time period $T+\Delta T$ then :-
\begin{qoptions}
\opt{A}{$\dfrac{\Delta T}T=\dfrac32\dfrac{\Delta r}r$}
\opt{B}{$\dfrac{\Delta T}T=\dfrac23\dfrac{\Delta r}r$}
\opt{C}{$\dfrac{\Delta T}T=\dfrac{\Delta r}r$}
\opt{D}{$\dfrac{\Delta T}T=-\dfrac{\Delta r}r$}
\end{qoptions}
\end{bankq}
\begin{bankq}[id=TMAS01-P1-PHY-S1-Q17, source={TMAS - Test 01 (09 Jul 2023)}, paper=1, qno=17, type=single, subtopic={Kepler's laws}, answer={C}, answer_src=key, difficulty=2, diff_parts={K1 S1 I0 C1 F0}, diff_note={concepts: elliptical orbit | steps: speed largest at perihelion}]
A planet revolves around the sun in an elliptical orbit. The linear speed of the planet will be maximum at
\bfig{87.9}{TMAS01_p1_phy_s1_q17_1.png}
\begin{qoptions}
\opt{A}{$D$}
\opt{B}{$B$}
\opt{C}{$A$}
\opt{D}{$C$}
\end{qoptions}
\end{bankq}
\begin{bankq}[id=TMAS01-P1-PHY-S2-Q06, source={TMAS - Test 01 (09 Jul 2023)}, paper=1, qno=6, type=numerical, subtopic={Field and potential}, answer={7}, answer_src=key, difficulty=6, diff_parts={K1 S2 I2 C1 F1}, diff_note={concepts: gravity inside a shell planet | steps: gravity at the surface; gravity at depth R; equate; solve the ratio}]
The density of core of a planet is $\rho_1$ and that of the outer shell is $\rho_2$. The radii of the core and outer shell of the planet are $R$ and $2R$ respectively. If acceleration due to gravity at the surface is half as at a depth $R$ then find $\dfrac{\rho_1}{\rho_2}$.
\end{bankq}
\begin{bankq}[id=TMAS01-P1-PHY-S2-Q07, source={TMAS - Test 01 (09 Jul 2023)}, paper=1, qno=7, type=numerical, subtopic={Field and potential}, answer={0}, answer_src=solved, working={Work per unit mass $=\vec E\cdot\vec d$ with $\hat d=(4,3)/5$: $(3\times4-4\times3)/5=0$}, difficulty=3, diff_parts={K1 S1 I0 C1 F1}, diff_note={concepts: work in a gravitational field | steps: field dot displacement}]
The gravitational field in a region is given by $\vec E=(3\hat i-4\hat j)$ N/kg. Find out the work done (in joule) in displacing a particle by 1 m, along the line $4y=3x+9$ :
\end{bankq}
\begin{bankq}[id=TMAS01-P2-PHY-S1-Q04, source={TMAS - Test 01 (09 Jul 2023)}, paper=2, qno=4, type=multi, subtopic={Kepler's laws}, answer={B,C,D}, answer_src=key, difficulty=9, diff_parts={K3 S2 I2 C2 F1}, diff_note={concepts: elliptical orbit, angular momentum, energy | steps: eccentricity from speeds; areal velocity; planet radius bound; energy change}]
A satellite revolves around a spherical planet in a stable elliptic orbit. Its speed at perigee is 2 km/s and its speed at apogee is 1 km/s. Its farthest distance from the planet's centre is 6000 km. Then which of the following statements are correct.
\begin{qoptions}
\opt{A}{The eccentricity of its orbit is 1/2.}
\opt{B}{The line joining the planet's center to the satellite sweeps out an area of $1.8\times10^{11}$ m$^2$ in each minute.}
\opt{C}{The planet's radius must be less than 3000 km.}
\opt{D}{The gain in its potential energy when it travels from the perigee to apogee is 3/4 times its kinetic energy at the perigee.}
\end{qoptions}
\end{bankq}
\begin{bankq}[id=TMAS01-P2-PHY-S1-Q05, source={TMAS - Test 01 (09 Jul 2023)}, paper=2, qno=5, type=multi, subtopic={Satellites}, answer={B,C,D}, answer_src=key, difficulty=6, diff_parts={K2 S2 I1 C1 F1}, diff_note={concepts: Kepler law, energy | steps: radius ratio; energy comparison; kinetic energy comparison}]
Two satellites of same mass of a planet in circular orbits have periods of revolution 32 days and 256 days. If the radius of the orbit of the first is $R$, then the
\begin{qoptions}
\opt{A}{Radius of the orbit of the second is 8$R$}
\opt{B}{Radius of the orbit of the second is 4$R$}
\opt{C}{Total mechanical energy of the second is greater than that of the first.}
\opt{D}{Kinetic energy of the second is less than that of the first.}
\end{qoptions}
\end{bankq}
\begin{bankpassage}[id=TMAS01-P2-PHY-S2-G11-13, source={TMAS - Test 01 (09 Jul 2023)}, questions={11-13}]
A particle is taken to a distance $r(>R)$ from centre of the earth. $R$ is radius of the earth. It is given velocity $V$ which is perpendicular to $\vec r$. With the given values of $V$ in column I you have to match the values of total energy of particle in column II and the resultant path of particle in column III. Here `$G$' is the universal gravitational constant and `$M$' is the mass of the earth.\par\smallskip
{\footnotesize\setlength{\tabcolsep}{1pt}\begin{tabular}{|l|l|l|l|l|l|}
\hline
 & \shortstack{Column-I\\(Velocity)} & & \shortstack{Column-II\\(Total Energy)} & & \shortstack{Column-III\\(Path)}\\
\hline
(I) & $V=\sqrt{\dfrac{GM}r}$ & (i) & Negative & (P) & Elliptical\\
(II) & $V=\sqrt{\dfrac{2GM}r}$ & (ii) & Positive & (Q) & Parabolic\\
(III) & $V>\sqrt{\dfrac{2GM}r}$ & (iii) & Zero & (R) & Hyperbolic\\
(IV) & $\sqrt{\dfrac{GM}r}<V<\sqrt{\dfrac{2GM}r}$ & (iv) & Infinite & (S) & Circular\\
\hline
\end{tabular}}
\end{bankpassage}
\begin{bankq}[id=TMAS01-P2-PHY-S2-Q11, source={TMAS - Test 01 (09 Jul 2023)}, paper=2, qno=11, type=single, subtopic={Satellites}, answer={A}, answer_src=key, difficulty=4, diff_parts={K1 S2 I1 C1 F0}, diff_note={concepts: orbit type from speed | steps: energy sign; circular, elliptical, parabolic, hyperbolic; match}, group=TMAS01-P2-PHY-S2-G11-13]
Choose the correct option :
\begin{qoptions}
\opt{A}{(I)-(i)-(S)}
\opt{B}{(II)-(iv)-(R)}
\opt{C}{(III)-(ii)-(Q)}
\opt{D}{(IV)-(iii)-(P)}
\end{qoptions}
\end{bankq}
\begin{bankq}[id=TMAS01-P2-PHY-S2-Q12, source={TMAS - Test 01 (09 Jul 2023)}, paper=2, qno=12, type=single, subtopic={Satellites}, answer={B}, answer_src=key, difficulty=4, diff_parts={K1 S2 I1 C1 F0}, diff_note={concepts: orbit type from speed | steps: energy sign; orbit shape; match}, group=TMAS01-P2-PHY-S2-G11-13]
Choose the correct option :
\begin{qoptions}
\opt{A}{(I)-(i)-(R)}
\opt{B}{(II)-(iii)-(Q)}
\opt{C}{(III)-(ii)-(S)}
\opt{D}{(IV)-(iv)-(P)}
\end{qoptions}
\end{bankq}
\begin{bankq}[id=TMAS01-P2-PHY-S2-Q13, source={TMAS - Test 01 (09 Jul 2023)}, paper=2, qno=13, type=single, subtopic={Satellites}, answer={A}, answer_src=key, difficulty=4, diff_parts={K1 S2 I1 C1 F0}, diff_note={concepts: orbit type from speed | steps: energy sign; orbit shape; match}, group=TMAS01-P2-PHY-S2-G11-13]
Choose the correct option :
\begin{qoptions}
\opt{A}{(III)-(ii)-(R)}
\opt{B}{(II)-(i)-(S)}
\opt{C}{(I)-(iii)-(P)}
\opt{D}{(IV)-(iv)-(Q)}
\end{qoptions}
\end{bankq}
\begin{bankq}[id=TMAS01-P2-PHY-S3-Q02, source={TMAS - Test 01 (09 Jul 2023)}, paper=2, qno=2, type=numerical, subtopic={Satellites}, answer={2}, answer_src=key, difficulty=8, diff_parts={K1 S3 I2 C2 F1}, diff_note={concepts: energy to form a satellite from a planet | steps: binding energy before; masses and orbit; energy after; difference; x}]
The energy required to dislodge 25\% mass of a planet of mass $M$ and radius $R$, so that is can form a satellite orbiting at a radius of 1.5 $R$ from the centre of the planet, is given by $E=\dfrac{0.21}x\dfrac{GM^2}R$. What is the nearest integer value of $x$? Assume that both the bodies are of spherical shape. [$G$ -- universal gravitational constant. Take $4^{1/3}=1.6$ and $\left(\dfrac43\right)^{1/3}=1.1$]
\end{bankq}
\begin{bankq}[id=TMAS01-P2-PHY-S3-Q03, source={TMAS - Test 01 (09 Jul 2023)}, paper=2, qno=3, type=numerical, subtopic={Field and potential}, answer={4}, answer_src=key, difficulty=6, diff_parts={K1 S2 I2 C1 F1}, diff_note={concepts: orbit in the field of two masses | steps: force components along the radius; centripetal equation; n}]
A particle of mass $m$ is moving in circular motion of radius $r$ with constant speed $v$ in influence of two fixed particles of mass $M$. Plane of circular motion is perpendicular to line joining fixed particles as shown. Speed of particle is given as $v=\sqrt{\dfrac{GM}{nr}}$ find value of $n$ ($a=r\sqrt3$)
\bfig{112.4}{TMAS01_p2_phy_s3_q03_1.png}
\end{bankq}
\begin{bankq}[id=TMAS01-P2-PHY-S3-Q04, source={TMAS - Test 01 (09 Jul 2023)}, paper=2, qno=4, type=numerical, subtopic={Satellites}, answer={3}, answer_src=key, difficulty=8, diff_parts={K2 S2 I2 C2 F1}, diff_note={concepts: ring around a planet, tension | steps: force on an element; centripetal part; tension from the element balance; n}]
A ring of mass $m$ and radius $3R$ is rotating with constant angular speed $\omega$ around a planet of mass $M$ and radius $R$. Center of ring and planet concide with each other. Tension in the ring is given as $T=\dfrac{GMm}{3n\pi R^2}$. Find value of $n$ $\left(\omega=\sqrt{\dfrac{GM}{9R^3}}\right)$.
\bfig{107.3}{TMAS01_p2_phy_s3_q04_1.png}
\end{bankq}

\section*{TNAS papers}
\begin{bankpassage}[id=TNAS02-P1-PHY-S3-G5-6, source={TNAS - Test 02 (07 May 2023)}, questions={5-6}]
An isosceles right angle tunnel ABO is made inside the earth from surface of earth to its centre as shown in the figure. The distance between the points A and B and that between the points B and O is $\dfrac{R}{\sqrt2}$, where $R$ is the radius of earth. The surface of tunnel is assumed to be frictionless. A small particle of mass $m_0$ is released from point A. The collision between the particle and the tunnel is assumed to be perfectly inelastic. Take mass of earth is $M$. Radius of earth $R=6400$ km and $g=10$ m/s$^2$. Now answer the following Questions.
\bfig{76.8}{TNAS02_p1_phy_s3_para5_6_1.png}
\end{bankpassage}
\begin{bankq}[id=TNAS02-P1-PHY-S3-Q05, source={TNAS - Test 02 (07 May 2023)}, paper=1, qno=5, type=numerical, subtopic={Field and potential}, answer={2513}, answer_src=solved, working={Motion along AB is SHM about B with $\omega=\sqrt{g/R}$ and takes a quarter period; the inelastic turn at B (perpendicular tunnels) stops the particle, then BO is again a quarter period. Total $=\pi\sqrt{R/g}=800\pi\approx2513$ s}, difficulty=8, diff_parts={K1 S3 I2 C2 F1}, diff_note={concepts: motion in a tunnel inside the earth | steps: force along each tunnel; SHM period; speed at the bend; inelastic stop; second leg}, group=TNAS02-P1-PHY-S3-G5-6]
The time taken by the particle to move in the tunnel from A to O (in sec).
\end{bankq}
\begin{bankq}[id=TNAS02-P1-PHY-S3-Q06, source={TNAS - Test 02 (07 May 2023)}, paper=1, qno=6, type=numerical, subtopic={Field and potential}, answer={0.25}, answer_src=key, difficulty=4, diff_parts={K1 S1 I1 C1 F1}, diff_note={concepts: work by gravity in the earth | steps: potential inside a uniform sphere; difference}, group=TNAS02-P1-PHY-S3-G5-6]
Work done by gravitational field of earth on the particle to reach from B to O is $x\dfrac{Gm_0M}{R}$. Find $x$.
\end{bankq}
\begin{bankq}[id=TNAS02-P2-PHY-S1-Q04, source={TNAS - Test 02 (07 May 2023)}, paper=2, qno=4, type=multi, subtopic={Satellites}, answer={A,C}, answer_src=key, difficulty=6, diff_parts={K1 S2 I2 C1 F1}, diff_note={concepts: binary stars from velocity graph | steps: period and amplitudes; orbit radii; mass relation; check each}]
Two stars are undergoing circular motion about their common centre of mass. Motion observed by an observes which is very far and in the plane of motion is SHM. Velocity versus phase rotated by both A and B is as shown. $R_1=8\times10^7$ km. (Take $G=\dfrac{20}{3}\times10^{-11}$ SI)
\bfig{105.1}{TNAS02_p2_phy_s1_q04_1.png}
\bfig{78.9}{TNAS02_p2_phy_s1_q04_2.png}
\begin{qoptions}
\opt{A}{Distance between both stars is $10^{11}$ m.}
\opt{B}{Distance between both stars is $10^8$ m.}
\opt{C}{Total mass of both stars is $\dfrac38\times10^{25}$ kg.}
\opt{D}{Total mass of both stars is $\dfrac78\times10^{25}$ kg.}
\end{qoptions}
\end{bankq}
\begin{bankq}[id=TNAS03-P1-PHY-S1-Q05, source={TNAS - Test 03 (28 May 2023)}, paper=1, qno=5, type=single, subtopic={Field and potential}, answer={D}, answer_src=key, difficulty=2, diff_parts={K1 S1 I0 C1 F0}, diff_note={concepts: apparent weightlessness | steps: centripetal condition; period}]
The time period of rotation of the earth is 24 hrs. If the earth increases its speed of rotation, at nearly what time period will the weight of body at the equator become zero ? Given $R=6400$ km and $g=10\ \mathrm{ms}^{-2}$
\begin{qoptions}
\opt{A}{$0.8$ hr}
\opt{B}{$1$ hr}
\opt{C}{$72$ min}
\opt{D}{$84$ min}
\end{qoptions}
\end{bankq}
\begin{bankq}[id=TNAS03-P1-PHY-S1-Q16, source={TNAS - Test 03 (28 May 2023)}, paper=1, qno=16, type=single, subtopic={Field and potential}, answer={A}, answer_src=key, difficulty=5, diff_parts={K1 S2 I2 C1 F0}, diff_note={concepts: field of a spherical shell | steps: mass enclosed; field formula; coefficient}]
Gravitational field at a distance $2R$ from the centre of a spherical shell of inner radius $R$, outer radius $3R$ and density $d$ is equal to $k(G.\pi.R.d)$, where $G$ is universal gravitational constant and $k$ is
\begin{qoptions}
\opt{A}{$\dfrac73$}
\opt{B}{$\dfrac53$}
\opt{C}{$3$}
\opt{D}{$\dfrac43$}
\end{qoptions}
\end{bankq}
\begin{bankq}[id=TNAS03-P1-PHY-S1-Q20, source={TNAS - Test 03 (28 May 2023)}, paper=1, qno=20, type=single, subtopic={Field and potential}, answer={A}, answer_src=key, difficulty=5, diff_parts={K1 S2 I2 C1 F0}, diff_note={concepts: potential of an annular disc | steps: potential of a ring; integrate over the annulus; mass density; value at P}]
A thin uniform annular disc (sec figure) of mass $M$ has outer radius $4R$ and inner radius $3R$. The work required to take a unit mass from point $P$ on its axis to infinity is
\bfig{91.6}{TNAS03_p1_phy_s1_q20_1.png}
\begin{qoptions}
\opt{A}{$\dfrac{2GM}{7R}\left(4\sqrt2-5\right)$}
\opt{B}{$-\dfrac{2GM}{7R}\left(4\sqrt2-5\right)$}
\opt{C}{$\dfrac{GM}{4R}$}
\opt{D}{$\dfrac{2GM}{5R}\left(\sqrt2-1\right)$}
\end{qoptions}
\end{bankq}
\begin{bankq}[id=TNAS03-P1-PHY-S2-Q01, source={TNAS - Test 03 (28 May 2023)}, paper=1, qno=1, type=numerical, subtopic={Kepler's laws}, answer={48}, answer_src=key, difficulty=6, diff_parts={K1 S2 I2 C1 F1}, diff_note={concepts: fall of a stopped planet | steps: degenerate ellipse; Kepler third law; fall time}]
If a planet were suddenly stopped in its orbit supposed to be circular, it would fall into the sun in a time $T\left(\dfrac{\sqrt2}{n}\right)$, where $T$ is the time period of revolution. Find $6n$.
\end{bankq}
\begin{bankq}[id=TNAS03-P1-PHY-S2-Q10, source={TNAS - Test 03 (28 May 2023)}, paper=1, qno=10, type=numerical, subtopic={Satellites}, answer={16}, answer_src=key, difficulty=4, diff_parts={K1 S1 I1 C1 F1}, diff_note={concepts: energy of a satellite | steps: orbital speed; kinetic energy in terms of T}]
In a satellite the time of revolution is $T$. Then K.E. is proportional to $T^n$. Find $(-24\ n)$
\end{bankq}
\begin{bankq}[id=TNAS03-P2-PHY-S3-Q06, source={TNAS - Test 03 (28 May 2023)}, paper=2, qno=6, type=numerical, subtopic={Satellites}, answer={452.5}, answer_src=key, difficulty=5, diff_parts={K1 S1 I2 C1 F1}, diff_note={concepts: orbit of a projectile | steps: energy and angular momentum; speed at the start}]
A projectile projected tangentially from the surface of a planet of radius $R$ (64 Km) from a point A, follows the indicated trajectory during its motion. If it is at a height of $3R$ at the farthest point of its trajectory, then the velocity of projection at A, $v_0$ is given by :- (Acceleration due to gravity on surface $g=2$ m/s$^2$)
\bfig{66.1}{TNAS03_p2_phy_s3_q06_1.png}
\end{bankq}
\begin{bankq}[id=TNAS03-P2-PHY-S3-Q07, source={TNAS - Test 03 (28 May 2023)}, paper=2, qno=7, type=numerical, subtopic={Satellites}, answer={2}, answer_src=key, difficulty=6, diff_parts={K1 S2 I2 C1 F1}, diff_note={concepts: bullet fired upward from a planet | steps: height from the field condition; energy; escape speed ratio}]
A bullet is fired vertically upwards with velocity $v$ from the surface of a spherical planet. When it reaches its maximum height, its acceleration due to the planet's gravity is $\dfrac14$th of its value of the surface of the planet. If the escape velocity form the planet is $v_{esc}=v\sqrt N$, then the value of $N$ is (ignore energy loss due to atmosphere)
\end{bankq}
\begin{bankq}[id=TNAS05-P1-PHY-S2-Q03, source={TNAS - Test 05 (09 Jul 2023)}, paper=1, qno=3, type=numerical, subtopic={Field and potential}, answer={32}, answer_src=key, difficulty=9, diff_parts={K2 S3 I2 C2 F1}, diff_note={concepts: particle grazing a planet, angular momentum and energy | steps: angular momentum conservation; energy conservation; grazing condition; solve for the speed; match the form}]
A particle of mass $m$ is thrown towards a stationary planet with speed $v$. So that it grazes the planet as shown in figure. Find $x$ if $v=\sqrt{\dfrac{xGM}{105R}}$.
\bfig{155.8}{TNAS05_p1_phy_s2_q03_1.png}
\end{bankq}
\begin{bankq}[id=TNAS05-P1-PHY-S2-Q05, source={TNAS - Test 05 (09 Jul 2023)}, paper=1, qno=5, type=numerical, subtopic={Escape velocity}, answer={2632}, answer_src=key, difficulty=8, diff_parts={K1 S3 I2 C2 F1}, diff_note={concepts: time to rise in the field of the earth with escape velocity | steps: escape condition; velocity as a function of radius; integrate dt; limits; evaluate; round}]
A body is projected vertically upwards from the surface of earth with a velocity just sufficient enough to carry it to infinity. The time taken by it to reach height $h=3R$ is \dots\dots S. (In nearest integer) (Take : $g=10$ m/s$^2$, Radius of earth : $R=6400$ km and $\sqrt2=1.41$)
\end{bankq}

\section*{Star Review (Other) papers}
\begin{bankq}[id=STAR03-P1-PHY-S1-Q02, source={Star Review (Other) - Test 03 (04 Dec 2023)}, paper=1, qno=2, type=single, subtopic={Field and potential}, answer={C}, answer_src=key, difficulty=4, diff_parts={K1 S2 I1 C1 F0}, diff_note={concepts: variation of gravitational acceleration with height | steps: expression for g at height; first order expansion; fractional change}]
The fractional change in the value of free--fall acceleration $g$ for a particle when it is lifted from the surface to an elevation $h$ ($h\ll R$) is
\begin{qoptions}
\opt{A}{$\dfrac hR$}
\opt{B}{$\dfrac{2h}{R}$}
\opt{C}{$-\dfrac{2h}{R}$}
\opt{D}{$-\dfrac hR$}
\end{qoptions}
\end{bankq}

\section*{Internal Test papers}
\begin{bankq}[id=INT01-P1-PHY-S1-Q01, source={Internal Test - Test 01 (05 Aug 2020)}, paper=1, qno=1, type=single, subtopic={Kepler's laws}, answer={A}, answer_src=key, difficulty=2, diff_parts={K1 S2 I0 C0 F0}, diff_note={concepts: orbit under a central electrostatic force | steps: which quantities are conserved; angular momentum conservation; speed varies along the ellipse}]
Under the influence of the coulomb field of charge $+Q$ which is fixed, a charge $-q$ is moving around it in an elliptical orbit. Find out correct statements
\begin{qoptions}
\opt{A}{The angular momentum of the charge $-q$ is constant.}
\opt{B}{The linear momentum of the charge $-q$ is constant.}
\opt{C}{The angular velocity of the charge $-q$ is constant.}
\opt{D}{The linear speed of the charge $-q$ is constant.}
\end{qoptions}
\end{bankq}
\begin{bankq}[id=INT01-P1-PHY-S3-Q15, source={Internal Test - Test 01 (05 Aug 2020)}, paper=1, qno=15, type=numerical, subtopic={Escape velocity}, answer={1.33}, answer_src=key, difficulty=6, diff_parts={K1 S2 I2 C1 F1}, diff_note={concepts: bullet fired upward and the escape speed | steps: height from the gravity ratio; energy conservation; escape velocity relation; N}]
A bullet is fired vertically upward with velocity ``$V$'' from the surface of a spherical planet. When it reaches its maximum height, its acceleration due to the planet's gravity is $\dfrac1{16}$th of its value at the surface of the planet. If the escape velocity from the planet surface is $V_{esc}=V\sqrt N$, then the value of $N$ is (ignore atmosphere).
\end{bankq}
\begin{bankq}[id=INT01-P1-PHY-S3-Q18, source={Internal Test - Test 01 (05 Aug 2020)}, paper=1, qno=18, type=numerical, subtopic={Field and potential}, answer={4}, answer_src=key, difficulty=6, diff_parts={K1 S2 I2 C1 F1}, diff_note={concepts: kinetic energy of two bodies falling together under gravity | steps: energy conservation from rest at infinity; total kinetic energy; coefficient}]
Two point objects of masses $m$ and $4m$ are at rest at an infinite seperation. A slight push is given to one of them and they move towards each other under mutual gravitational attraction. If $G$ is the universal gravitational constant. Then at a seperation ``$r$'' total K.E of objects is given by $\dfrac{xGm^2}{r}$ find $x$?
\end{bankq}
\begin{bankq}[id=INT02-P1-PHY-S1-Q05, source={Internal Test - Test 02 (07 Jul 2021)}, paper=1, qno=5, type=single, subtopic={Field and potential}, answer={D}, answer_src=key, difficulty=3, diff_parts={K1 S2 I1 C0 F0}, diff_note={concepts: two bodies attracting from rest at infinity | steps: energy and momentum conservation; centre of mass; kinetic energy at separation r}]
Two objects of masses $m$ and $4m$ are at rest at infinite separation. They move towards each other under mutual gravitational attraction. Then, at a separation $r$, which of the following is true?
\begin{qoptions}
\opt{A}{The total kinetic energy of the system is not zero.}
\opt{B}{The force between them is not zero.}
\opt{C}{The centre of mass of the system is at rest.}
\opt{D}{All the above are true.}
\end{qoptions}
\end{bankq}
\begin{bankq}[id=INT02-P1-PHY-S2-Q09, source={Internal Test - Test 02 (07 Jul 2021)}, paper=1, qno=9, type=multi, subtopic={Satellites}, answer={A,D}, answer_src=key, difficulty=5, diff_parts={K1 S2 I1 C1 F1}, diff_note={concepts: angular momentum of a satellite | steps: orbital speed; angular momentum; relation with the surface gravity; direction}]
A satellite of mass $m$ is moving in a circular orbit of radius $r$ around a planet of mass $M$
\begin{qoptions}
\opt{A}{The magnitude of angular momentum with respect to the centre of the orbit is $m\sqrt{GMr}$, where $G$ is the gravitational constant.}
\opt{B}{The magnitude of the angular momentum is $mR\sqrt{2gr}$ where $g$ is the acceleration due to gravity on the surface of the planet.}
\opt{C}{The direction of angular momentum is parallel to the plane of the orbit.}
\opt{D}{The direction of angular momentum is perpendicular to the plane of the orbit.}
\end{qoptions}
\end{bankq}
\begin{bankq}[id=INT02-P1-PHY-S3-Q13, source={Internal Test - Test 02 (07 Jul 2021)}, paper=1, qno=13, type=numerical, subtopic={Escape velocity}, answer={0.575}, answer_src=key, difficulty=5, diff_parts={K1 S2 I1 C1 F1}, diff_note={concepts: speed at the surface when dropped from a height | steps: energy conservation; speed at the surface; ratio to escape velocity}]
A body is dropped from a height $h$ equal to $\dfrac R2$, where $R$ is the radius of the earth. If it will hit the surface of the earth with a speed $v=av_e$, where $v_e$ is the escape velocity from the surface of the earth. Find $a$
\end{bankq}
\begin{bankq}[id=INT02-P1-PHY-S3-Q14, source={Internal Test - Test 02 (07 Jul 2021)}, paper=1, qno=14, type=numerical, subtopic={Field and potential}, answer={0.5}, answer_src=key, difficulty=5, diff_parts={K1 S2 I1 C1 F1}, diff_note={concepts: weight change in a tunnel | steps: variation of g with depth; fractional decrease; percentage}]
Assuming the earth to be a sphere of uniform mass density, find the percentage decrease in the weight of a body when it is taken to the end of a tunnel $32$ km below the surface of the earth. Radius of earth $=6400$ km
\end{bankq}

\section*{Nurture (General) papers}
\begin{bankq}[id=NUR01-P1-PHY-S1-Q03, source={Nurture (General) - Test 01 (12 Jun 2021)}, paper=1, qno=3, type=multi, subtopic={Escape velocity}, answer={B,C,D}, answer_src=solved, working={g is proportional to r cubed inside; escape speed equals that of a uniform planet; energy is proportional to mass}, difficulty=6, diff_parts={K2 S2 I1 C1 F1}, diff_note={concepts: planet with density proportional to r squared, escape velocity | steps: mass inside radius r; g inside; escape speed; compare with uniform planet}]
A spherical planet of radius $R$ has spherically symmetrical distribution of mass density, varying as square of the distance from the centre, from zero at centre to maximum value $\rho_0$ at its surface.
\begin{qoptions}
\opt{A}{The value of escape velocity of a mass $m$ at the surface of planet is $\sqrt{\dfrac{4\pi G\rho_0R^2}{5}}$.}
\opt{B}{The value of acceleration due to gravity '$g$' varies inside the planet as cube of the distance from centre.}
\opt{C}{The value of escape velocity from surface is same as the escape velocity from another planet of same total mass \& radius but having uniform mass density.}
\opt{D}{The energy required to impart escape velocity to particles of masses '$m$' \& '$2m$', at the surface of planet, will be in ratio 1:2.}
\end{qoptions}
\end{bankq}
\begin{bankq}[id=NUR04-P1-PHY-S1-Q03, source={Nurture (General) - Test 04 (11 Feb 2024)}, paper=1, qno=3, type=single, subtopic={Field and potential}, answer={D}, answer_src=solved, working={Momentum conservation and energy GMm/d split the kinetic energy; v=sqrt(2G M\_other\^{}2/(d(m+M)))}, difficulty=5, diff_parts={K2 S2 I1 C1 F0}, diff_note={concepts: speeds of two masses falling together under gravity, field and potential | steps: momentum conservation; energy conservation; solve for speeds}]
Two particles of masses m and M are initially at rest at an infinite distance apart. They move towards each other and gain speeds due to gravitational attraction. Find their speeds when the separation between the masses becomes equal to d are
\begin{qoptions}
\opt{A}{$V_1=\sqrt{\dfrac{GM^2}{2d(m+M)}};V_2=\sqrt{\dfrac{Gm^2}{2d(m+M)}}$}
\opt{B}{$V_1=\sqrt{\dfrac{Gm^2}{d(m+M)}};V_2=\sqrt{\dfrac{GM^2}{d(m+M)}}$}
\opt{C}{$V_1=\sqrt{\dfrac{2Gm^2}{d(m+M)}};V_2=\sqrt{\dfrac{2GM^2}{d(m+M)}}$}
\opt{D}{$V_1=\sqrt{\dfrac{2GM^2}{d(m+M)}};V_2=\sqrt{\dfrac{2Gm^2}{d(m+M)}}$}
\end{qoptions}
\end{bankq}
\begin{bankq}[id=NUR04-P1-PHY-S3-Q15, source={Nurture (General) - Test 04 (11 Feb 2024)}, paper=1, qno=15, type=numerical, subtopic={Satellites}, answer={625}, answer_src=solved, working={F=6r\^{}2=mv\^{}2/r gives KE=375 J, U=250 J, E=625 J}, difficulty=6, diff_parts={K2 S2 I1 C1 F1}, diff_note={concepts: particle in a cubic potential in circular orbit, satellites | steps: force from the potential; centripetal condition; kinetic energy; total energy}]
A particle of mass 2kg is moving under the influence of a force which always acts towards the center and whose potential energy is given by $U(r)=2r^3j$. If the body is moving in a circular orbit of radius 5m, then find its energy in J.
\end{bankq}
\end{document}
